\chapter{CA 19-9 (Cancer Antigen 19-9)}

\textbf{Interpretive caution:} CA 19-9 levels are not a diagnostic test for pancreatic cancer and should never be used to establish the diagnosis of pancreatic cancer. Up to 10--15\% of the general population does not express the CA 19-9 epitope (Lewis a-b- genotype), so their CA 19-9 levels are unmeasurable even in the presence of advanced disease. CA 19-9 is also elevated in many benign conditions (chronic pancreatitis, cholangitis, jaundice, gallstones, cirrhosis, and other malignancies such as cholangiocarcinoma, gastric cancer, and colorectal cancer). CA 19-9 should be interpreted in the clinical context of symptoms, imaging, and histologic findings.

\section*{What Is This Test?}
The CA 19-9 test measures the concentration of carbohydrate antigen 19-9 (CA 19-9, sialyl-Lewis a, sLe$^\text{a}$) in the blood. CA 19-9 is a sialylated Lewis blood group antigen (a type of mucin). It is expressed on the surface of normal pancreatic ductal cells and biliary epithelial cells and in high amounts on the surface of pancreatic cancer cells (especially ductal adenocarcinoma). It is the most widely used serum tumor marker for pancreatic cancer, but it is NOT a screening test and is not specific for pancreatic cancer. It is used for (1) supporting the diagnosis in a symptomatic patient with a pancreatic mass and clinical suspicion; (2) monitoring response to chemotherapy or surgery; (3) detecting recurrence after curative-intent resection. CA 19-9 is elevated (>37 U/mL) in 65--85\% of patients with resectable pancreatic cancer and in >90\% of patients with metastatic disease. However, because of its limitations (false negatives in Lewis-negative individuals, false positives in benign biliary/pancreatic disease), it is not used for screening or for establishing a diagnosis.

\section*{Alternative Names}
\begin{itemize}
  \item Carbohydrate Antigen 19-9
  \item CA 19-9
  \item Cancer Antigen 19-9
  \item sLe$^\text{a}$
  \item Sialyl Lewis a Antigen
  \item Pancreatic Cancer Marker
  \item GI Tumor Marker CA 19-9
\end{itemize}
\section*{Principle}
CA 19-9 is measured by immunoassay (chemiluminescent microparticle immunoassay, CMIA, or electrochemiluminescent immunoassay, ECLIA). The assay uses a monoclonal antibody (1116-NS19-9) raised against a colorectal cancer cell line, which recognizes the sialylated Lewis a (sLe$^\text{a}$) antigen. CA 19-9 is a sialylated blood group determinant (Lewis a epitope with sialic acid), so its expression depends on the Lewis blood group genotype: approximately 90\% of the population has the Lewis a+b+ or a-b+ genotype and can express CA 19-9; the remaining 5--10\% are Lewis a-b- and cannot synthesize CA 19-9 at all (these individuals have undetectable CA 19-9 even with advanced pancreatic cancer). CA 19-9 is elevated in pancreatic ductal adenocarcinoma because tumor cells express sLe$^\text{a}$ and the antigen is shed into the bloodstream. CA 19-9 is also elevated in cholangiocarcinoma, gallbladder cancer, gastric cancer, colorectal cancer, and hepatocellular carcinoma, as well as in benign biliary obstruction, acute and chronic pancreatitis, jaundice, cholangitis, gallstones, cirrhosis, and hepatic abscesses. CA 19-9 is cleared by the liver and excreted in the bile, so hepatobiliary obstruction or impaired hepatic function increases CA 19-9 levels.

\section*{History}
CA 19-9 was discovered in 1979 by Hilary Koprowski and colleagues at the Wistar Institute. They immunized mice with a colorectal cancer cell line (SW1116) and developed the monoclonal antibody 1116-NS19-9, which reacted with a sialylated Lewis blood group antigen \cite{koprowski1979}. The biochemical characterization of the antigen as sialyl-Lewis a was published by Magnani and colleagues in 1981 \cite{magnani1981}. It was initially evaluated as a marker for colorectal cancer, but it showed greater utility for pancreatic cancer. It became the most widely used tumor marker for pancreatic cancer. Major limitations were recognized: false negatives in Lewis-negative individuals, false positives in benign biliary and pancreatic disease, and low sensitivity for early-stage disease. Contemporary guidelines (including the 2021 NCCN and ASCO guidelines) recommend CA 19-9 for prognosis, response monitoring, and recurrence surveillance in pancreatic cancer, but not for screening \cite{tempero2021}.

\section*{Why This Test Is Ordered}
\begin{itemize}
  \item Evaluation of a suspected pancreatic mass: CA 19-9 is measured in a patient with weight loss, jaundice, or a newly discovered pancreatic mass on imaging. A markedly elevated CA 19-9 (>100--200 U/mL) in this setting supports the diagnosis of pancreatic ductal adenocarcinoma, although biopsy is still required for definitive diagnosis
  \item Monitoring response to systemic chemotherapy in metastatic pancreatic cancer: a decline in CA 19-9 (>20--50\% from baseline) during treatment is associated with better outcomes, and a rise suggests progression
  \item Post-resection surveillance for recurrence: CA 19-9 is measured every 3--6 months for 2 years after curative-intent pancreatectomy. A rise above 37 U/mL (or doubling from the postoperative nadir) suggests recurrence
  \item Assessment of resectability and prognosis: CA 19-9 >500 U/mL at presentation is associated with occult metastases and poor prognosis; CA 19-9 normalization after neoadjuvant chemotherapy predicts a favorable outcome after surgery
  \item Not for screening: CA 19-9 is not recommended for screening asymptomatic individuals because of low sensitivity for early-stage disease and low specificity (elevated in benign pancreaticobiliary disease and other cancers)
\end{itemize}

\section*{Primary vs Secondary Test}
CA 19-9 is a secondary/supplementary test. It never establishes a diagnosis by itself. It supports the diagnosis of pancreatic cancer when combined with symptoms, imaging (CT/MRI/EUS), and histology (biopsy or FNA). Its primary role is monitoring response to therapy and detecting recurrence.

\section*{How This Test Is Performed}
A healthcare professional draws blood from a vein into a serum separator tube (gold-top). The sample is centrifuged and analyzed by immunoassay. Results are typically available within 1--2 hours. Serial measurements should be performed with the same assay and laboratory for reliable trending.

\section*{What Preparation Is Needed}
Fasting (8--12 hours) is often recommended, as eating may transiently elevate CA 19-9. However, non-fasting values are still clinically interpretable in most settings. Document current medications and any recent endoscopic procedures (ERCP can transiently elevate CA 19-9).

\section*{Contraindications and Precautions}
\begin{itemize}
  \item CA 19-9 is NOT a diagnostic test. It cannot distinguish benign from malignant pancreaticobiliary disease. Benign biliary obstruction (choledocholithiasis, cholangitis), acute and chronic pancreatitis, and cirrhosis cause CA 19-9 elevations that frequently overlap with the levels seen in early pancreatic cancer
  \item CA 19-9 is undetectable or very low in 5--10\% of the population (Lewis a-b- genotype). In these individuals, CA 19-9 is unhelpful even in the presence of advanced pancreatic cancer. If CA 19-9 is unexpectedly low/undetectable, consider Lewis antigen testing
  \item CA 19-9 is elevated in other GI malignancies: cholangiocarcinoma (biliary marker), gastric cancer, colorectal cancer, hepatocellular carcinoma, and also in lung, breast, and gynecologic cancers
  \item Jaundice falsely elevates CA 19-9: biliary obstruction raises CA 19-9 levels independent of tumor burden. In a jaundiced patient, the CA 19-9 should be rechecked after biliary decompression (stent or surgery) for accurate interpretation
  \item CA 19-9 is not useful for screening asymptomatic patients (poor sensitivity for early disease) and is not a substitute for tissue diagnosis (EUS-FNA or surgical biopsy)
  \item CA 19-9 is elevated in benign conditions such as endometriosis, ovarian cysts, uterine fibroids, and in pregnancy (up to 20 U/mL). A mildly elevated CA 19-9 (37--100 U/mL) is common in healthy individuals and does not indicate cancer
  \item Serial CA 19-9 should be measured with the same assay and laboratory: inter-assay variability can cause spurious trends
\end{itemize}
\section*{Risks}
Standard venipuncture risks: mild pain, bruising, or hematoma at the puncture site.

\section*{What Happens After the Test}
Results are available within 1--2 hours. CA 19-9 is interpreted with imaging and clinical findings. In a patient with a pancreatic mass, an elevated CA 19-9 supports malignancy but does not replace tissue biopsy. If CA 19-9 was ordered for treatment monitoring, the trend is more informative than any single value. If CA 19-9 is elevated in a patient with jaundice, it should be rechecked after biliary decompression.

\section*{Understanding Results}

\subsection*{Below Normal}
\begin{itemize}
  \item \textbf{CA 19-9 <37 U/mL:} Normal reference range. This does NOT exclude pancreatic cancer: up to 15--20\% of patients with pancreatic cancer have normal CA 19-9 (including all Lewis-negative patients and some patients with small or early-stage tumors)
  \item \textbf{Undetectable or very low CA 19-9:} Suspect the Lewis a-b- genotype. In such patients, CA 19-9 is unmeasurable even in advanced disease and is not useful for monitoring
  \item \textbf{Normalization of CA 19-9 after surgery:} A return to <37 U/mL after curative-intent resection is a favorable prognostic sign and is the baseline for future surveillance
  \item \textbf{Low CA 19-9 during chemotherapy:} A decline from a previously elevated baseline indicates tumor response (Rustin-like criteria). A >20--50\% reduction correlates with improved survival
\end{itemize}

\subsection*{Above Normal}
\begin{itemize}
  \item \textbf{CA 19-9 37--100 U/mL:} Mildly elevated. May be seen in benign conditions (chronic pancreatitis, gallstones, cholangitis, cirrhosis) and in early pancreatic cancer. Interpretation requires clinical context; serial measurement may be helpful
  \item \textbf{CA 19-9 >100--200 U/mL:} Moderately elevated. In a patient with a pancreatic mass, this level strongly supports pancreatic ductal adenocarcinoma. In the absence of a pancreatic mass, it should prompt evaluation for other causes (cholangiocarcinoma, gastric cancer, colorectal cancer, hepatobiliary disease)
  \item \textbf{CA 19-9 >500 U/mL:} Markedly elevated. In pancreatic cancer, this predicts a high likelihood of occult metastatic disease and poor resectability. Preoperative CA 19-9 >500 U/mL is associated with a lower rate of R0 resection and worse survival
  \item \textbf{CA 19-9 rising during chemotherapy:} Suggests disease progression. A doubling or a rise above 1.5--2x the nadir is commonly used as a trigger for restaging imaging
  \item \textbf{CA 19-9 rising after curative-intent resection:} First sign of recurrence (may precede imaging by 1--3 months). A value >37 U/mL (or a doubling from the postoperative nadir) on two consecutive samples indicates biochemical recurrence and warrants restaging imaging
\end{itemize}

\section*{Normal Results}
\begin{itemize}
  \item \textbf{Adults:} 0--37 U/mL (reference range varies by assay and laboratory)
  \item \textbf{Healthy individuals:} Most have CA 19-9 <20 U/mL; up to 5\% of healthy adults have CA 19-9 >37 U/mL
  \item \textbf{Lewis a-b- individuals:} CA 19-9 is undetectable or very low (<2 U/mL) regardless of disease
\end{itemize}

\section*{Follow-up Tests}
\begin{itemize}
  \item \textbf{CT abdomen/pelvis with contrast (pancreatic protocol):} First-line imaging for a suspected pancreatic mass. A hypodense, hypovascular pancreatic mass with upstream ductal dilatation is characteristic of pancreatic adenocarcinoma
  \item \textbf{EUS with fine-needle aspiration (EUS-FNA):} Gold standard for tissue diagnosis of pancreatic masses. Cytology plus CA 19-9 levels from the aspirate (not serum) are sometimes used, but serum CA 19-9 remains the standard blood marker
  \item \textbf{Endoscopic retrograde cholangiopancreatography (ERCP):} For biliary decompression and brushings in jaundiced patients with biliary obstruction; also for tissue diagnosis when EUS is unavailable
  \item \textbf{Liver function tests and bilirubin:} Jaundice and cholestasis elevate CA 19-9. CA 19-9 should be interpreted with concurrent bilirubin: values are unreliable in the setting of obstruction
  \item \textbf{CEA (carcinoembryonic antigen):} Often elevated in pancreatic cancer and other GI malignancies; used in combination with CA 19-9 in some clinical scenarios
  \item \textbf{Repeat CA 19-9 after biliary decompression:} If jaundice was present at baseline, CA 19-9 should be remeasured after stent placement or surgical drainage. A persistent elevation after relief of obstruction is more specific for malignancy
\end{itemize}

\section*{References}
\bibliographystyle{unsrtnat}
\bibliography{references}

\clearpage
\section*{Regulatory Status and Authorization Date}
This chapter describes a test category, not a specific commercial product. FDA, CDSCO, EU IVDR, and MHRA approval or registration dates therefore cannot be assigned without the assay manufacturer, product name, instrument, and intended use. For laboratory-developed tests, an individual agency approval date may not exist. See the Preface for the interpretation of regulatory status.

\clearpage
